\answer{ex:17:1}{$0.2$~W so với $81$~\textmu W; luồng thô cần không quá $50$~pJ/bit.}
\answer{ex:17:2}{Khi $c=0.1$: $5.6$, $8.7$, $9.2$~bit/s, giới hạn $9.3$~bit/s; khi $c=0$ và $N=1000$: $65$~bit/s.}
\answer{ex:17:3}{$B=\log_2N+P\log_2P+(1-P)\log_2\frac{1-P}{N-1}=4.87$, $4.37$, $3.29$~bit mỗi ký tự, tức $7.3$, $6.6$, $4.9$~bit/s; để đạt $10$~bit/s cần $2.3$~ký tự mỗi giây.}
\answer{ex:17:4}{Phương sai sai số $2b/(Nm^2T)+\sigma_\xi^2$ mỗi thành phần; hai phần đóng góp bằng nhau khi $N\approx90$; giới hạn $\tfrac12\log_2(1+0.25/0.04)\approx1.4$~bit mỗi thành phần mỗi cửa sổ.}
\answer{ex:17:5}{Cặp hội tụ khi $\eta_bd^2+\eta_db^2<2$, tức khoảng $\eta<1$: với $0.8$ hội tụ, với $1.1$ dao động ổn định chu kỳ $2$, với $1.5$ không tuần hoàn, với $2$ phân kỳ; cần tách tốc độ học của hai bên.}
\answer{ex:17:6}{Ngưỡng $\approx4.4\sigma$; $102$~kbit/s, $0.10$~mW so với $205$~mW cho luồng thô, nhỏ hơn $2000$ lần; chỉ $5.7\cdot10^{-5}$ số giá trị đếm bị bão hòa (hơn $3$ xung).}
\answer{ex:17:7}{Bài mở. Khoảng $46$~bit/s khi $\text{SNR}\approx0.9$ và khoảng $330$~bit/s với nhiễu độc lập. Nếu cản trở là nhiễu giống tín hiệu, thông tin rút ra sẽ bão hòa khi số kênh tăng; nếu là do chưa hiểu mã, nó sẽ tăng với bộ giải mã tốt hơn (chẳng hạn bộ giải mã dùng thời điểm xung).}
